Measuring the surface area of an object from a CT scan is harder than it looks, because the scan is a
grid of small cubes (voxels) and the true surface passes between them. Counting the exposed faces of
the voxels gives an area about 50\% too large at every resolution. Marching cubes, which fits a
triangle mesh to the data, is closer, but its error depends on how curved the object is and how noisy
the scan is. We built 2,400 synthetic shapes from seven families whose true surface area is known to
better than 0.01\%, and turned them into 14,400 CT-like volumes at $32^3$ and $64^3$ voxels with three
noise levels. We then trained models to predict, for each object, how much the marching-cubes area
should be corrected: a Random Forest on ten simple shape measurements, a 3D convolutional neural
network (CNN) and a PointNet. On test shapes that none of the models had seen, marching cubes has an
average error of 4.78\% ($32^3$) and 3.52\% ($64^3$). The best corrections bring this down to 0.68\%
(3D CNN, $32^3$) and 0.18\% (Random Forest, $64^3$), and stay below 2\% on two shape families left out of
training altogether. The simple Random Forest does as well as the CNN overall, the CNN is better only
for objects smaller than 16 voxels, and no method can recover structures thinner than about two voxels.
